\documentclass[10pt,a4paper]{amsart}
\usepackage[margin=19mm]{geometry}
\usepackage{amsmath,amssymb,mathtools}
\usepackage[scr=rsfs,cal=euler]{mathalfa}
\usepackage{xfrac}
\usepackage[unicode,hidelinks,bookmarks=false]{hyperref}
\newcommand{\ie}{i.e.\ }
\newcommand{\cf}{cf.\ }
\newcommand{\resp}{respectively\ }
\newcommand{\Subsection}{\subsection}
\providecommand{\nobreakdash}{\nobreak\hbox{-}}
\setlength{\emergencystretch}{2em}
\multlinegap0pt
\usepackage{amssymb}
\usepackage{tikz}
\usepackage{mathtools}
\newtheorem{theorem}{Theorem}
\newcommand{\e}{\mathbf{o}}
\newcommand{\rmN}{\mathbf{N}}
\newcommand{\varN}{\mathcal{N}}
\newcommand{\varP}{\mathcal{P}}
\newcommand{\x}{\mathbf{x}}
\title{A Finite Automaton Approach to Combinatorial Games}
\author{Kai Liang}
\date{Source: arXiv:2608.13273v2 (CC BY 4.0)}
\begin{document}
\maketitle

\noindent\emph{Source and licence.} A Finite Automaton Approach to Combinatorial Games. Version arXiv:2608.13273v2 (CC BY 4.0), distributed by arXiv under the Creative Commons Attribution 4.0 International licence, http://creativecommons.org/licenses/by/4.0/, as recorded in the arXiv licence metadata for this exact version and shown on the abstract page https://arxiv.org/abs/2608.13273v2. Full licence notice: \texttt{licenses/arxiv-2608.13273-CC-BY-4.0.txt}.\par\smallskip

\noindent\textit{Editorial scope.} Counted unit: the article's theorem line\_dfa (source0.tex lines 1343--1349). The article's complete original proof of it is delivered with it (source0.tex lines 1351--1407, one \texttt{proof} environment of 57 lines). No mathematics is written, repaired or re-derived: the statement and the proof are the article's own lines, in the article's own notation, and no retained line is substituted or re-flowed.  No further source-local context is needed: every cross-reference of every retained line resolves inside the two delivered spans.  Nothing was omitted from any retained span, so the package carries no omission record.  No retained line contains a citation, so the wrapper carries no bibliography and no bibliography entry.  No retained line contains a figure environment, an image-inclusion command or drawing code, so no image is delivered and the package declares no asset dependency.  Editorial formatting only, in addition: an AMS \texttt{amsart} preamble that restates the article's own declarations and macros used by the retained lines, namely the environments \texttt{theorem} on separate counters, together with the standard packages those lines need.  The retained environments print in this document as Theorem 1 in order of appearance, whereas the article prints the same items with the numbering of its own full document.  The complete native source of the article as distributed in the arXiv e-print for this exact version is bundled unchanged as \texttt{sources/207577/source0.tex}.
\begin{theorem}\label{line_dfa}
    For any FA-line game, any seed $\gamma\in\varGamma$, and any label $\theta_0 \in\varTheta$, the set of character parts of all $\varN$-positions with that seed and label,
    \[
        W = \{w \mid w\in\varSigma^*,\ \gamma w\theta_0 \in \varN\},
    \]
    is a regular language over $\varSigma$. The same holds for $\varP$-positions.
\end{theorem}

\begin{proof}
    Unlike the usual direction, the DFA constructed in this proof reads characters from right to left.
    To obtain an equivalent left-to-right DFA, one simply reverses all state transitions and then determinises.

    We take the set of all seeds $\varGamma$ as the state set of the DFA, and construct the transition mapping according to the recursive outcome relations,
    so that starting from the initial state $\gamma$, it reads characters one by one from right to left, and after all characters are read, the outcome is given according to the token label.
    The key point that guarantees finiteness of the state set is that the legal move lengths are bounded:
    since a seed $\gamma$ contains $s_{\max}$ positions (the length of the seed), after a new position is recorded by $\gamma$, the number of positions becomes $s_{\max}+1$, but the rightmost position is unreachable by any legal move and can therefore be deleted, keeping the number of positions bounded by $s_{\max}$.

    If the position is $\theta_0\gamma$, i.e., there are no characters between the token and the seed, then any legal move of the token enters the seed directly, and the outcome can be computed immediately:
    \[
        o(\theta_0\gamma) =
        \begin{cases}
            \varN, & \exists s\in S,\ \gamma_s(\varDelta(\bar{\gamma}_{1}\ldots\bar{\gamma}_{s-1}\e ,\theta_0)) = \rmN; \\
            \varP, & \text{otherwise}.
        \end{cases}
    \]
    This gives the final states of the DFA: $\gamma$ is a final state if and only if $\theta_0\gamma$ has the corresponding outcome.

    For the transition mapping, if an obstacle $\x$ is encountered, we simply delete the rightmost column and add a column of obstacles $\x$ on the left:
    \[
        \x\cdot
        \begin{bmatrix}
            \gamma_1(\theta_1) & \ldots & \gamma_{s_{\max}}(\theta_1) \\
            \gamma_1(\theta_2) & \ldots & \gamma_{s_{\max}}(\theta_2) \\
            \vdots & & \vdots \\
            \gamma_1(\theta_{|\varTheta|}) & \ldots & \gamma_{s_{\max}}(\theta_{|\varTheta|}) \\
        \end{bmatrix}
        =
        \begin{bmatrix}
            \x & \gamma_1(\theta_1) & \ldots & \gamma_{s_{\max}-1}(\theta_1) \\
            \x & \gamma_1(\theta_2) & \ldots & \gamma_{s_{\max}-1}(\theta_2) \\
            \vdots & \vdots & & \vdots \\
            \x & \gamma_1(\theta_{|\varTheta|}) & \ldots & \gamma_{s_{\max}-1}(\theta_{|\varTheta|}) \\
        \end{bmatrix}.
    \]

    If an empty cell $\e$ is encountered, we delete the rightmost column and add a column on the left whose entries are computed according to the legal moves for each label:
    \[
        \e\cdot
        \begin{bmatrix}
            \gamma_1(\theta_1) & \ldots & \gamma_{s_{\max}}(\theta_1) \\
            \gamma_1(\theta_2) & \ldots & \gamma_{s_{\max}}(\theta_2) \\
            \vdots & & \vdots \\
            \gamma_1(\theta_{|\varTheta|}) & \ldots & \gamma_{s_{\max}}(\theta_{|\varTheta|}) \\
        \end{bmatrix}
        =
        \begin{bmatrix}
            \e_{o(\theta_1 \gamma)} & \gamma_1(\theta_1) & \ldots & \gamma_{s_{\max}-1}(\theta_1) \\
            \e_{o(\theta_2 \gamma)} & \gamma_1(\theta_2) & \ldots & \gamma_{s_{\max}-1}(\theta_2) \\
            \vdots & \vdots & & \vdots \\
            \e_{o(\theta_{|\varTheta|} \gamma)} & \gamma_1(\theta_{|\varTheta|}) & \ldots & \gamma_{s_{\max}-1}(\theta_{|\varTheta|}) \\
        \end{bmatrix}.
    \]

    This gives the complete structure of a DFA satisfying the requirements.
\end{proof}

\end{document}
